\subsection{Robots!}\label{robots}

Robots! The word is loaded with both promise and peril. We dream of
robots that do all tedious labor, freeing humanity from it, as well as
of robots that might take over and kill us all(fiction seems to favor
the latter).

I also believe robots can be transformative, although I think we should
look at much of the hype from today's ``tech'' companies with many
grains of salt. Self driving cars and autonomous battle robots have
mostly turned out to be worthless hype machines useful for making
Silicon Valley hucksters rich and not for much else.

Here I will look at some of the robots I think we should build with
Trash Magic which can make a better world for caring for one another and
having adventures, which is what this book is all about.

\subsection{A Rumble of Robots}\label{a-rumble-of-robots}

The collective noun for robots is ``a rumble of robots.'' I'm not sure
where I heard this, I think on of my friends may have made it up, but
it's so perfect it's too good not to use. So I want to talk about
rumbles of robots. In particular the difference between robots used for
consumption and for production.

Amazon is in the process of building robot based infrastructure for
delivery. This is fundamentally a consumption driven project. The main
initial figure of merit in the growth of their network will be coverage:
the more potential consumers are covered, the better. This will mean
that it is optimal for robots to be as far as possible from other
robots. But how does this picture change for production?

Rumbles of robots are very common on the production side of things.
Those who produce cars and computers and the like often have rumbles of
robots, with humans just as technicians who run the machines. Much like
a cow hand or shepherd, I think there should be a name for those who
herd rumbles of robots: rumbler. So the trash wizard is also a rumbler.
And the trash wizard stick is like the shepherd's crook: a tablet that
drives a IoT network that consists of your rumble of robots.

That is what seizing the means of production is really all about. It's
not about seizing an existing factory, which will be based on existing
methods, or about building a primitive system that can't compete. It's
about building rumbles of robots which can reproduce themselves by
harvesting free materials to make more, and then rumbling them around to
build what else you need.

Key elements of the trash wizards robot rumble are mobility and
versatility. They will run off of locally harvested energy, and be
programmed to gather energy as needed as well as materials. They should
scale in that the robots you need for a 10 bot rumble are not so
different from a single roninbot or a 1000 bot rumble. They should be
able to reproduce from found materials and forage for those materials
with some simple guidance from the rumbler. That is the plan.

\subsection{The Tripod and Multitool}\label{the-tripod-and-multitool}

The basis for all Trash Magic fabrication will be the tripod. a tripod
is exactly what it sounds like: three sticks, joined at the top with
wire or tape or rope, as well as with guy wires at the bottom if needed,
with the bases of the three sticks planted in the ground for stability.
A rock hangs from the center to keep it steady.

Each of the three sticks shall be of Skeletron. Each shall be the size
of the Standard Trash Magic Staff. A practitioner of Trash Magic shall
have the power to deploy this at any time, at any scale and with any
tool. Full xyz motion is possible with wires that are fed through holes
or wheels on the sticks, and can go to stepper motors for drive and/or
wheels on potentiometers. Each wire connects to an arduino that runs a
three terminal pulsed transport measuring protocol. The ADC signals are
amplified and converted with some math to an audio signal, visual
display, and or vibrotactile interface.

This probe can be heated with the wires, can do local electrochemical
work and measurement, measure nonlinearities of all kinds, and can thus
weld, plate, edit, implant ions, plastic weld and plastic re-shape,
signals also displayed in a GUI in a smart phone if needed but can also
not be, can have no ``smart'' computer at all, just dumb computers.
Sensing can be capacitive, inductive, resistive, frequency dependent or
not, nonlinear or not, gated or not. Many different sizes and weights
can be used, gear ratios can be changed on the fly with various
interchangeable wheels. Everything is modular.

Thermal regulation of processes can be carried out by moving the
position of a suspended cauldron in 3 dimensions relative to a fire at
the bottom of the tripod, using a microprocessor and feedback on a
thermometer to regulate the temperature to a fixed point.

In a cluster of Trash Magic industrial production there can be many
tripods, all in a network, connected by plastic plumbing, electrical and
data connections, and zip lines that move materials from one to another,
setting up ad hoc nodes in an assembly line that can be grown and shrunk
on the fly.

Tripods can also support roof and wall materials, making shelters that
people can sleep and work in. Tripods can be trees that already exist.
Tripods can also be completely submerged for doing intricate work on the
bottom of a body of water. In zero gravity another set of sticks are
needed, and another cable and springs. The materials for a pair of
tripods should also be able to become the skeleton of a simple boat,
which can then deploy mini tripods to the water below it.

Tripods can be fractal, with a small tripod lowered from the tool head
of a large tripod. Smallest tripods should have Angstrom resolution and
atomically sharp tips, biggest ones are hundreds of meters tall.

Elements:

3 skeletron sticks

wire or rope for lashing

knife and drill and pliers to work the stuff together

trash magic philosophy engine style motor(one sided four pole stepper,
water proof). all made from plastic bottles, HDPE milk bottles and HDPE
caps from soda bottles, coils unwound and re wound from big old AC
motors and transformers. Transistors and SoC devices are salvaged and
hacked from internet routers. Micro controllers are salvaged from
thermostats and the like. Op amps are salvaged from audio electronics of
various kinds. motors designed to run as ``dumb'' steppers or DC
brushless fast motors with feedback using pickup coil(s).

IOT device that has enough outputs to drive all the motors, e.g.~the Pi
zero or equivalent SoC

smart phone or tablet to control the IOT thing remotely--if there is VR
or AR, it happens on this machine so that it's off the shelf and tracks
whatever the tech assholes come up with as it develops.

sensor wheels that can pick up position of a passing structural wire as
it goes past as the work moves, along with arduino to read them all out
which connects to the IOT platform and out to the tablet, this allows
full xyz input by the user or recording of position while a user works
the machine by hand

LED illumination pointed at work space from the Trash Magic sticks,
these have RGB capability so they can signal various types of data

speakers and amplifier that can play sound from the tablet, IOT device
or direct amplification to get audio feedback on the signal

arduino that does the three terminal generic electrical measurement and
drive with pulse width modulation and various ADC pickups, this connects
to the raspberry pi, which then puts out light, sound, and vibration to
feed back electrical data to the user

various heated tools and electrical heaters which can be driven by the
suspending wires--need lots of thermometers for feedback, which can go
to arduinos that connect to the pi over serial or bluetooth

hydro electric generators with tethers so that the whole tripod can be
run by the tide or stream nearby, close but not necessarily under the
tripod(although they could be!)

generic pumps for water and gas--two drive coils and a pickup coil, this
is also the vibrational driver, see below, same tool, can be either on
the moving piece or stationary while the work piece is moved.

vibrating tool which can vibrate a heavy rock-weighted tool that can,
over many hours, under moving water, carve arbitrary shapes into stone,
vibration also used for plastic welding with heat

soldering tool connected to Main Trash Magic Power

designs for many sizes and types of wheel so that torque can be
transformed up or down, all fabricatable from bottles.

LiPo batteries and bigger ad hoc electrochemical cells for storing more
energy, possibly a water ``tower'' where water with stored potential
energy can flow pressurized fluid through the system

high voltage generator for plasma work etc

\subsection{Octahedral Ball Drone}\label{octahedral-ball-drone}

The octahedral ball drone is a octahedron made of three intersecting
sticks, with a flexible joint. Each of the 6 ends has a 2 degree of
freedom +/- pair of control coils and magnets with some significant
amount of effective length change. Drivers can use the natural dynamics
of the flexing rods to make efficient rolling motion. I imagine making
the first one based on a LiPo battery and Raspberry Pi, with wooden
sticks held with a rubber central ball that the sticks can get pushed
into. I imagine the sticks being a total of about 12 to 15 inches long
total, so that each pod is about 6-8 inches from the center, and can
flex by about 2-3 inches, or 20\% to 50\% of total. The amount of play
should be just barely more than enough to get the full sphere effect for
maximally efficient rolling. I imagine for the very first version all
coils are driven by dumb feedback and it's an art piece that rolls
around in all directions chaotically until the battery dies. Then all
the boards get tied together to a central Raspberry Pi, which gets onto
the net for control by a remote console. Then the code task gets solved
for the simplest possible rolling motion with directional bias. Then
some kind of simple chase algorithm is written so that it can track down
a target that emits some kind of signal. This demo could launch a whole
robotics project which can be a main support column for The Project.

\subsection{Wind Desert Drift Robots}\label{wind-desert-drift-robots}

Rolling robots with windmills, they roll, then gather wind electricity
into a capacitor, roll again, and repeat. They can go for hundreds of
miles with no intervention. The instinct to go a certain direction based
on navigating off of the sun is programmed into the physical hardware.
After some long time, maybe many years, the machine calls for help,
eventually someone finds it and follows the instructions for repair and
improvement. With generation after generation editing and helping the
thing exist, it can exist for hundreds of years, slowly cleaning up
wasted sacrafice zones of the old capitalist world.

There are so many machines to build like this! Machines that comb the
ocean for contaminants, using waves go get energy to move around and
sort and grab stuff, potentially floating around for years before being
found based on a data stream that pulses out periodically, and
eventually another type of robot tending robot can grab it, extract the
materials it's gathered, and bring it to a floating factory robot
rumble. This kind of robot is important for the ecosystem of the jungle
city in the ocean-inundated coastal post apocalypse.

Free robots like this are a rational response to the fact that the
existing system has created sacrafice zones. These sacrafice zones have
negative economic value in the old system, making them freely available
to be absorbed into the anarchist industrial infrastructure. This is
key: in order to avoid getting crushed by the forces of the old system
too early our movement must exist in the fringes of the current system,
where the old ways have created land of negative value. The very fact
that land can have negative value, that this is a concept that people
accept, should be yet another red flag that assignment of numerical
values to real human values is a morally bankrupt act.

This should always be the goal of free technology if it wants to grow
exponentially without a lot of resistance: the input must be things
deemed of ``negative value'' by the old system. Unlike most projects in
capitalism which constantly drain everyone involved more and more over
time, creating generation after generation of institutional burnout.

\subsection{Cyborgs}\label{cyborgs}

birds, fungi, marine mammals, cephalopods , trees, rodents,
trees\ldots{}

\begin{figure}[htbp]
\centering
\includegraphics{images/birdcyborg.png}
\caption{Bird Cyborg}
\end{figure}

\begin{figure}[htbp]
\centering
\includegraphics{images/squidcyborg.png}
\caption{Squid Cyborg}
\end{figure}

\subsection{Building an Ent}\label{building-an-ent}

The fractal mater reactor should be alive. Trees, bushes, grass, etc.
can grow all around it, with roots going into various fractal channels
which provide nutreince. These liquid spaces can have various animals
and fungi and microorganisms, creating a whole ecosystem. Imagine an
island built up of such mater, the size of a small building, covered
with trees. Ambient energy is used to slowly build up and discharge
electrical energy to operate philosophy engines which slowly walk the
whole thing across the landscape. With little or even no human
intervention, this limbering living giant might spend decades crawling
up and down hills scouring for junk cars, which it turns into a ever
growing robot rumble that it can give away to any passing humans for
free at any time. Building this kind of thing in the ocean can be
incredibly powerful. Whole floating islands filled with fractal reactor
technology can wander the high seas, with the humans all underwater in
bubbles to ride out storms, picking up storm energy and sea junk, and
building a every larger floating city deep out in the ocean. This
aquatic fractal techno city could exist even in a dead world of violent
storms and acid oceans and extreme heat. That's part of why I hate the
liberal vision of fighting global warming by trying to make everything
into Denmark. It won't actually work, and then you still have a
breakable society. Sure our society is destroying itself, but so what?
Let's build a million new ones, not just one. And let's expect a future
not of competent bureaucrats carefully tabulating the giant World
Spreadsheet so everyone can live the life of a middle class urban Sweede
but a world where ice ages and supervolcanoes and nuclear wars and
devastating earthquakes all happen and where we roll with it and have
adventures.

Start by building and documenting some actual robots, make a rumble,
rumble the rumble with the trash wizard stick.

\begin{enumerate}
\def\labelenumi{\arabic{enumi}.}
\tightlist
\item
  hopper dumb robot
\item
  hopper with a brain using raspberry pi for browser to pi to board to
  motor robot control
\item
  roller robot
\item
  rolling ball robot
\item
  hopper rumble
\item
  make a mobile robot that can use a hot tool to rework thermo plastics
\item
  floating robot that can re work thermo plastic(loop back to these on
  fractal reactor chapter)
\end{enumerate}

\subsection{Flying Robots}\label{flying-robots}

Everyone is in such a hurry! Most aerial drones for personal use
today(2016) are designed to move very fast for very short periods of
time. Generally with four propellors pointed straight up, they can take
off fast anywhere, go in all directions fast, dodge fast moving
obstacles, and often only last a few minutes. If broken, they have
numerous small parts which can be very hard to fix.

Quad copter personal drones are great capitalist technology: they break
easily, cost a lot, do very little, need constant upgrades, and are
mostly ``useful'' for entertaining the techno-priesthood and annoying
everyone else. Not surprisingly I see much that can be improved here.

The first way I would set about making drones less useless is by making
them float instead of fly with propellors. Given that they're both small
and don't have living cargo, I would say the arguments against hydrogen
for lift are mostly obsolete.

\begin{figure}[htbp]
\centering
\includegraphics{images/dirigible.png}
\caption{Hydrogen Drone}
\end{figure}

How should motors work for soaring drones? First of all, if the thing is
large enough it can float on the circular current patterns in the upper
atmosphere, holding position with no mechanical work done. But what
about motors for guidance? These motors should be electrostatic, powered
by extremely high voltage giant balloon capacitors which are the main
body of the soaring drone. Using two very light polymers in very thin
sheets with opposite positions on the turboelectric series, it should be
possible when far from the ground to generate \emph{extremely} high
voltages very easily using the mechanical energy source of the rotating
air currents. Electrostatic motors can then run off these, also built
from thin polymer sheets with thin metallization. No magnets and no
copper coils! It's nutty to use the magnetic field for high altitude low
power low speed motors, they should all be based on electric fields,
because it's easy to get megavolts up there.

\subsection{How Robots Reproduce}\label{how-robots-reproduce}

Not on their own! With help. Robots can always ask for help, and it is
our task as their designers and creators to build the information into
them in the form of works of art that makes it obvious how to repair and
extend the robot. A robot should also be constructed in such a way that
it is its own means of production: the components of the robot can be
used as a machine to build more robots like it. This will require human
effort, but both the physical tools and the information required to
learn the skills to duplicate the machine are built into the machine and
obvious to find and use. Modern technology is designed to scare you away
from modifying it or interacting with it in any deep way. We seek to
build machines that do the opposite: invite the user to get more deeply
involved, building more, documenting that process, and extending the
technology themselves for others to use.

I will illustrate this with an example. One of the simplest robots will
move itself around looking for energy, then when it finds some(generally
a fast moving water body like a creek) it will turn itself into both a
power plant and a chemical plant, storing energy and chemicals extracted
from the water(targeting human industrial waste of various kinds). This
will involve a computer, some motors and some sensors. Other machines
will be involved in large scale computer fabrication as outlined in
other sections of this work.

\subsection{Littoral Robots}\label{littoral-robots}

Add: littoral robot rumbles which can use tides and river currents to
generate electricity to propel themselves upstream. Can be amphimious,
use water to charge but land to move, can move with hopping, jumping,
walking, rolling, and slithering. Littoral trash cleanup robots are
fundable, can make a huge difference in cleanup of a waterway, and also
give us free source material for more building.

\subsection{Hopping, rolling, walking, running, jumping,
driving}\label{hopping-rolling-walking-running-jumping-driving}

\subsection{Robots with different times scales, centuries of work, or
hours of
lifetime}\label{robots-with-different-times-scales-centuries-of-work-or-hours-of-lifetime}

Something that I think needs to be investigated in robot design is time
scale. Capitalists like a certain time scale--the shorter the better.
But without capitalism and its obsession with short term growth and
profits we can set times scales on hundreds of years or even long in
some cases. Suppose an area of land is contaminated with plutonium or
some other radioactive heavy metal. It might be there for many thousands
of years, making the land uninhabitable. Thousands of years, but not
forever, and plutonium has uses even in a peaceful society without
rules. Why not clean it up?

Perhaps the robots that clean plutonium will grow their own biofuel to
get energy from the sun and slowly pick their way across the land,
working with cyborg worms and fungi to dig up the atoms and move them
together and out of the water table. How many processes of atomic or
molecular transport open up when we allow a process to take thousands of
years? Many. I'm sure capitalists already use the term ``geological
engineering'' but I would say that to truly apply that term, you should
be carrying out a technical/artistic endeavor which takes place on a
geological time scale. That means it has to be \emph{very} easy for
future people to understand, maintain and repair. It also has to
anticipate future geological changes, including catastrophic ones like a
volcano that destroyed all life on earth for a billion years. And it
should have time horizons that stretch well into the 10's of millions of
years. What's your hurry! If we were not all hounded by debt to
capitalists we could take time to really work on hard things like
plutonium cleanup one atom at a time.

\subsection{rumble jacks}\label{rumble-jacks}

Imagine the following drone/robot: magnets on duct tape with 2 drive
coils each are on the vertices of a octahedron made from a star of
sticks with the brain and energy storage in the middle. With the right
kind of intelligence, all the magnets can move in a coordinated way to
roll the whole ball like a jack in the game jacks. Each jack is aware of
the position of its neighbors, and together they make a rumble of robots
that act like a herd of sheep. Hence ``rumbleJack'' for the rumble of
jacks like robots.

This technology can be used for all sorts of long slow land cleaning
processes. Rather than try to maximize battery life, they will use
capacitors to store energy, and recharge the capacitors from ambient
energy. For a rumble of jacks in the prairie, the obvious source of
power is the wind. Ideally, the wind will be used to create energy which
will immediately go into directed propulsion. This might be slow since
it depends on gusts, but it can go on forever, so slowness becomes ok.
This is technology that you would deploy to spend 1000 years cleaning up
a sacrifice zone, where you want no outside energy or materials to be
needed at all and for the rumble to keep doing its work for hundreds of
years. Also, obviously, clearing of mine fields is a immediate
application. A rumble of tire-sized octahedra could potentially roll
themselves at 10's of miles per hour, keeping up with a car or truck and
making it possible for the rumble to proceed in a mob ahead of a motor
vehicle, taking out IED's in real time. The rumble could end up in a
convoy geometry, stretched out over the length of the road, doing recon
ahead and tracking behind to see what's happening after a convoy passes.
In these applications it probably makes sense for the source of power to
be the trucks or cars in the human/freight convoy, with individuals in
the rumble cycling through the charging station and back out into the
rumble.

Going back to the plaines discussed above, this is a great tool for
agriculture. Even just gathering. A gathering rumble could go out and
gather roots and berries from the countryside in a quasi-cultivated
area. These roving balls could be picking up and dropping seeds as they
go, mapping where all the useful plants are, and also harvesting as they
go, taking wind sun and water as energy sources as needed, then spending
energy when it's available to do the work.

Another robot rumble I want to build that is closely related is the
slithering water robots. These use the usual magnet and coil arrangement
to create a slithering motion in buoyant objects, which can then
smoothly cut through the water. The fact that this has not been widely
deployed is totally insane: the same drive can be used in reverse to get
electrical power out of wave action. If the length of each robot is a
few wavelengths, the whole thing will be forced into a wave which can
create EMF as the magnets move, which can go into the storage
capacitors, then released to change slightly the nature of the
serpentine motion to direct the drone in a specific direction.

These can be incredibly powerful technology! The ocean can be a
fantastic source of raw materials for the trash wizards. Note that for
neutrally buoyant drones, this can serve to move them through the water
below the surface. One mode of operation might be to cruise a few meters
above the bottom of the ocean, scanning for stuff to salvage, then dive
and grab rocks to be negatively buoyant once a target is found. With
just barely negative buoyancy, the rumble can float just above the
target as they pick it apart. They then drop the weights, rise up,
inflate bags to float(everything is made from rubber, and reversible
air/vacuum/water pumps are in all things), and pull up and bring
material back to assembly centers, which can also be floating robot
rumble factories. With ocean currents and waves as an energy source, and
no hurry, these robots can work as slow as they have to, slowly making
more and more of themselves until they can have a global impact on ocean
cleanup.

The water based propulsion system also is very appealing for boats. I
want a boat that runs on wave action, wind, and tides, to grab energy as
it finds it, and then use it as needed to move toward a destination. I
can imagine this being just about kayak or canoe sized. I could also
imagine a freighter that is meters or even 10's of meters long. That
sounds small for a freighter, but imagine, again, that they're a huge
rumble that can be easily scaled up. This can be a freight swarm to move
materials across water.

One more thing I want to talk about relating to these technologies: the
serpentine gliding drone. This would be made out of very thin light
frame with large thin polymer wings to allow for gliding. Magnet-coil
drives are used to pulse the shape to optimize gliding, and to interact
with turbulence(linear wind won't work because we need relative motion
up in the air) to acquire energy,probably in the form of \emph{Very}
high voltages because that is most compatible with being high up.
Perhaps electric fields that exist due to natural weather can also be
used for various electrical things up high. A flying swarm that can
glide and gather energy and never have to land if the conditions are
right has many applications. Mapping to gain information on materials
for trash wizards to salvage is an obvious application. This type of
flying serpentine glider can also be used as hardware for data transfer
by flying flash memory. Donors like this are also a fantastic source of
much more granular weather and climate data than are now available. This
is useful in the long term for practical climate and weather science
studies to deal with mitigation during the coming storms of the next 200
years from climate change. But they also have use for short term weather
sensing tasks: tracking down storms that can be used as energy sources
for very large scale projects.

What about scaling these robots way up in size and weight for use inside
storms? One could imagine giant metal gliders in massive rumbles of 10's
of thousands or maybe even millions of units, all ripping around in a
storm could over the ocean. These generate giant hydrogen-filled blimps
which then gather in a huge rumble to go turn back into useful work near
a settlement or floating factory.

The Anthropocene is here. Like it or not, it's here. For the next 1000
years our planet is going to be dominated by the actions we choose to
take as a civilization. If we stay on the track we're on, the atmosphere
and oceans heat up, massive desertification destroys wet ecosystems
while rising oceans eat most of our cities, and the oceans become a
toxic waste dump that cannot sustain life. If we do nothing that is
clearly what will happen. Or something worse involving nuclear
holocaust. Given these alternatives, what difference does it make how
drastically we change things in the sea, air, and land? The opportunity
to simply not let civilization get big enough to destroy the world has
long passed us by now.

So is it so wrong to imagine the whole landscape filled with these
lumbering rumbles of rolling, slithering, hopping, and gliding robots?
Is it wrong to let them reproduce with human help, but with very little
labor-time, allowing groups of people to build endlessly expanding
rumble spheres around the world to create a world of total abundance? I
say that it is not wrong. Maybe if there were a way to go back it would
be a hard choice to do something that disturbs the balance of nature
like this, but there isn't.

This is what the trash wizard wants to make possible in the world:
endless streams of material and data moving through the physical world
with robots made from trash, which encompass our whole human
environment. Maybe not the whole world, but enough of it. A world of
abundance using the rumble sphere and value circles could exist outside
of the states and corporations. It does not need land, just someplace to
move to--it is all mobile by default. The trash wizards build the needed
expertise up and document it and teach it so that any group of people
can create this kind of culture anywhere, specific to their individual
cultural needs and the available resources in whatever geographical area
they're in.
